\FloatBarrier\section{Arquitecturas seleccionadas y justificación}

\subsection{ResNet-18 como referencia residual}
ResNet introduce conexiones que suman una transformación aprendida a una ruta de identidad. La formulación básica es $y=F(x)+x$ cuando las dimensiones son compatibles. Si cambian resolución o número de canales, una proyección lineal $W_s x$ ajusta la ruta de atajo. Esta estructura facilita la optimización de redes profundas, pero no significa que toda red más profunda sea mejor para cualquier dataset \cite{he2016}.

ResNet-18 utiliza bloques residuales básicos $3\times 3$ y constituye una referencia ampliamente documentada. En este proyecto interesa porque ofrece una arquitectura suficientemente expresiva sin recurrir de inmediato a una variante mayor. La cabeza original de ImageNet (512 entradas) se reemplaza por una capa lineal de siete salidas, resultando en 11\,180\,103 parámetros reales.

Como hipótesis, la referencia residual puede ofrecer una comparación clara entre capacidad y costo. No se presupone que domine en sensibilidad de melanoma ni en precisión global. La evidencia deberá provenir de la misma partición y receta utilizadas con los otros modelos.

\subsection{EfficientNet-B0 y el equilibrio de recursos}
EfficientNet estudia cómo escalar de forma coordinada la profundidad ($d=\alpha^\phi$), el ancho ($w=\beta^\phi$) y la resolución ($r=\gamma^\phi$) de una familia de redes mediante un coeficiente compuesto $\phi$. B0 es la referencia base de la familia y utiliza bloques invertidos eficientes (\textit{MBConv}) y mecanismos de atención de canales (\textit{Squeeze-and-Excitation}, SE) \cite{tan2019}.

El proyecto no implementa una búsqueda de nuevas arquitecturas ni compara B0 con toda la familia. Utiliza la implementación oficial de TorchVision con pesos \texttt{IMAGENET1K\_V1} y adapta la capa final de su clasificador (\texttt{1280} a \texttt{7} salidas, manteniendo \texttt{Dropout(0.2)}), sumando 4\,016\,515 parámetros totales. El artículo original informa ventajas de eficiencia bajo sus propias condiciones; esas cifras no se trasladan al proyecto como resultados locales sin medición directa.

\subsection{MobileNetV3-Small y la alternativa compacta}
MobileNetV3 combina búsqueda orientada al hardware, bloques invertidos con cuello de botella, convoluciones separables en profundidad, módulos SE y activaciones \textit{hard-swish} en capas profundas \cite{howard2019}. Su variante Small adaptada reemplaza la última capa lineal (\texttt{1024} a \texttt{7} salidas, tras la proyección intermedia $576\to 1024$ con \textit{Hardswish} y \textit{Dropout(0.2)}), sumando 1\,525\,031 parámetros.

La convolución separable divide el filtrado espacial por canal y la combinación puntual $1\times 1$. Esto reduce parámetros y operaciones respecto a una convolución convencional. Sin embargo, el acceso a memoria, el tamaño del lote y la implementación del operador influyen en la latencia real sobre GPU. El experimento debe medir tanto el desempeño diagnóstico como la velocidad real en la misma tarjeta gráfica.

\begin{table}[htbp]\centering\small
\setlength{\tabcolsep}{6pt}
\caption{Especificación técnica y parámetros exactos de las tres CNN adaptadas a siete clases (VRAM sintética con batch 2 a $224\times 224$).}
\begin{tabularx}{\textwidth}{l>{\raggedright\arraybackslash}Xlrr}\toprule
\textbf{Arquitectura} & \textbf{Mecanismo central} & \textbf{Pesos base} & \textbf{Parámetros} & \textbf{VRAM sintética} \\ \midrule
\textbf{ResNet-18} & Bloques residuales $3\times 3$ & \texttt{IMAGENET1K\_V1} & 11\,180\,103 & 147 MiB \\
\textbf{EfficientNet-B0} & MBConv + Squeeze-and-Excitation & \texttt{IMAGENET1K\_V1} & 4\,016\,515 & 208 MiB \\
\textbf{MobileNetV3-Small} & Depthwise separable + SE + h-swish & \texttt{IMAGENET1K\_V1} & 1\,525\,031 & 58 MiB \\
\bottomrule\end{tabularx}\end{table}


\begin{table}[htbp]\centering\footnotesize
\setlength{\tabcolsep}{4.5pt}
\caption{Comparación estructural de bloques, dimensiones de representación y cabezas adaptadas.}
\begin{tabularx}{\textwidth}{@{}>{\bfseries}lccc>{\raggedright\arraybackslash}X@{}}
\toprule
\textbf{Modelo} & \textbf{Activación} & \textbf{Atención SE} & \makecell{\textbf{Dim.\ entrada}\\\textbf{cabeza}} & \textbf{Capa sustituida en \texttt{models.py}} \\ \midrule
ResNet-18 & ReLU & No & 512 & \texttt{fc: Linear(512, 7)} \\
EfficientNet-B0 & SiLU (Swish) & Sí & 1\,280 & \texttt{classifier[-1]: Linear(1280, 7)} \\
MobileNetV3-Small & ReLU / h-swish & Sí & 1\,024 & \texttt{classifier[-1]: Linear(1024, 7)} \\
\bottomrule
\end{tabularx}
\end{table}

\figura{11_arquitecturas.png}{0.98}{Esquema estructural comparativo de las tres CNN adaptadas. Diagrama conceptual de etapas y dimensiones verificadas en \texttt{models.py}.}
\FloatBarrier

\subsection{Condiciones para una comparación justa}
Los tres modelos utilizarán el mismo orden de clases, split, resolución, política de aumento, pérdida, criterio de selección y presupuesto máximo de épocas. Se registran las versiones de los pesos ImageNet, porque ``preentrenado'' no identifica por sí solo una receta. La normalización de canales se mantiene compatible con los pesos, mientras que la geometría común mediante relleno constituye una decisión experimental documentada.

El estudio compara arquitecturas bajo una receta común. No pretende demostrar el óptimo absoluto de cada una, que requeriría búsquedas equivalentes de hiperparámetros y mayor presupuesto. VGG, ResNet-50 o ViT-B/16 no se descartan por ser incorrectos, sino porque no son necesarios para responder la pregunta inicial con cuatro GiB de VRAM. YOLO y U-Net corresponden a otras anotaciones y objetivos; no se incorporan como clasificadores por simple disponibilidad de código.
